\documentclass{amsart}
% For dealing with references we use the comment environment
\usepackage{verbatim}
\newenvironment{reference}{\comment}{\endcomment}
%\newenvironment{reference}{}{}
\newenvironment{slogan}{\comment}{\endcomment}
\newenvironment{history}{\comment}{\endcomment}

% For commutative diagrams we use Xy-pic
\usepackage[all]{xy}

% We use 2cell for 2-commutative diagrams.
\xyoption{2cell}
\UseAllTwocells

% We use multicol for the list of chapters between chapters
\usepackage{multicol}

% This is generally recommended for better output
\usepackage{lmodern}
\usepackage[T1]{fontenc}

% For cross-file-references

% Package for hypertext links:
\usepackage{hyperref}

% For any local file, say "hello.tex" you want to link to please

% Theorem environments.
%
\theoremstyle{plain}
\newtheorem{theorem}[subsection]{Theorem}
\newtheorem{proposition}[subsection]{Proposition}
\newtheorem{lemma}[subsection]{Lemma}

\theoremstyle{definition}
\newtheorem{definition}[subsection]{Definition}
\newtheorem{example}[subsection]{Example}
\newtheorem{exercise}[subsection]{Exercise}
\newtheorem{situation}[subsection]{Situation}

\theoremstyle{remark}
\newtheorem{remark}[subsection]{Remark}
\newtheorem{remarks}[subsection]{Remarks}

\numberwithin{equation}{subsection}

% Macros
%
\def\lim{\mathop{\mathrm{lim}}\nolimits}
\def\colim{\mathop{\mathrm{colim}}\nolimits}
\def\Spec{\mathop{\mathrm{Spec}}}
\def\Hom{\mathop{\mathrm{Hom}}\nolimits}
\def\Ext{\mathop{\mathrm{Ext}}\nolimits}
\def\SheafHom{\mathop{\mathcal{H}\!\mathit{om}}\nolimits}
\def\SheafExt{\mathop{\mathcal{E}\!\mathit{xt}}\nolimits}
\def\Sch{\mathit{Sch}}
\def\Mor{\mathop{\mathrm{Mor}}\nolimits}
\def\Ob{\mathop{\mathrm{Ob}}\nolimits}
\def\Sh{\mathop{\mathit{Sh}}\nolimits}
\def\NL{\mathop{N\!L}\nolimits}
\def\CH{\mathop{\mathrm{CH}}\nolimits}
\def\proetale{{pro\text{-}\acute{e}tale}}
\def\etale{{\acute{e}tale}}
\def\QCoh{\mathit{QCoh}}
\def\Ker{\mathop{\mathrm{Ker}}}
\def\Im{\mathop{\mathrm{Im}}}
\def\Coker{\mathop{\mathrm{Coker}}}
\def\Coim{\mathop{\mathrm{Coim}}}

% Boxtimes
%
\DeclareMathSymbol{\boxtimes}{\mathbin}{AMSa}{"02}

%
% Macros for moduli stacks/spaces
%
\def\QCohstack{\mathcal{QC}\!\mathit{oh}}
\def\Cohstack{\mathcal{C}\!\mathit{oh}}
\def\Spacesstack{\mathcal{S}\!\mathit{paces}}
\def\Quotfunctor{\mathrm{Quot}}
\def\Hilbfunctor{\mathrm{Hilb}}
\def\Curvesstack{\mathcal{C}\!\mathit{urves}}
\def\Polarizedstack{\mathcal{P}\!\mathit{olarized}}
\def\Complexesstack{\mathcal{C}\!\mathit{omplexes}}
% \Pic is the operator that assigns to X its picard group, usage \Pic(X)
% \Picardstack_{X/B} denotes the Picard stack of X over B
% \Picardfunctor_{X/B} denotes the Picard functor of X over B
\def\Pic{\mathop{\mathrm{Pic}}\nolimits}
\def\Picardstack{\mathcal{P}\!\mathit{ic}}
\def\Picardfunctor{\mathrm{Pic}}
\def\Deformationcategory{\mathcal{D}\!\mathit{ef}}

\setlength{\emergencystretch}{3em}
\begin{document}
\title[Stacks Project, Tag 0GTU]{567 --- Quasi-compact quasi-separated schemes admit proper finitely presented completely decomposed covers with ample families of line bundles}
\maketitle
\noindent Copyright (C) 2005--2025 Johan de Jong. Stacks Project authors. Source: \href{https://stacks.math.columbia.edu/tag/0GTU}{Tag 0GTU}.

\noindent Editorial difficulty note (not author text). Difficulty H2: The proof must preserve residue-field lifts while arranging an ample family on the covering scheme. Noetherian approximation reduces to finite-dimensional schemes, where an admissible blowup handles a dense open and dimension induction supplies a complementary proper cover with the same residue-field lifting property..

\noindent Editorial provenance note (not author text). History: excerpt from revision \texttt{a0444\allowbreak 6e57e\allowbreak c1fbc\allowbreak 252a8\allowbreak 71afc\allowbreak ec775\allowbreak 2fb28\allowbreak 07b14}, more-morphisms.tex, lines 23440--23483. Reference presentation was adapted; original-author context and core support blocks were retained or added. The mathematical wording is unchanged. Licensed under GNU FDL 1.2 or later; no invariant sections or cover texts. See ../licenses/COPYING. Context blocks reproduce author definitions and supporting results; they are not additional counted problems.

\begin{definition}
\label{definition-cd-morphism}
A morphism $f : X \to Y$ of schemes is said to be
{\it completely decomposed}\footnote{This may be nonstandard terminology.}
if for all points $y \in Y$ there
is a point $x \in X$ with $f(x) = y$ such that the field
extension $\kappa(x)/\kappa(y)$ is trivial.
A family of morphisms $\{f_i : X_i \to Y\}_{i \in I}$ of
schemes with fixed target is said to be {\it completely decomposed}
if $\coprod f_i : \coprod Y_i \to X$ is completely decomposed.
\end{definition}

\begin{proposition}
\label{proposition-envelope-with-resolution-property}
Let $X$ be a quasi-compact and quasi-separated scheme. There exists a
morphism $f : Y \to X$ which is of finite presentation, proper, and
completely decomposed (Definition \ref{definition-cd-morphism})
such that the scheme $Y$ has an ample family of invertible modules.
\end{proposition}

\begin{proof}
By Limits, Proposition \href{https://stacks.math.columbia.edu/tag/01ZA}{proposition approximate (Tag 01ZA)}
there exists an affine morphism $X \to X_0$ where $X_0$
is a scheme of finite type over $\mathbf{Z}$. Below we produce
a morphism $Y_0 \to X_0$ with all the desired properties.
Then setting $Y = X \times_{X_0} Y_0$ and $f$ equal to
the projection $f : Y \to X$ we conclude.
To see this observe that $f$ is of finite presentation
(Morphisms, Lemma \href{https://stacks.math.columbia.edu/tag/01TS}{lemma base change finite presentation (Tag 01TS)}),
$f$ is proper
(Morphisms, Lemma \href{https://stacks.math.columbia.edu/tag/01W4}{lemma base change proper (Tag 01W4)}),
$f$ is completely decomposed
(Lemma \href{https://stacks.math.columbia.edu/tag/0GTK}{lemma base change cd (Tag 0GTK)}). Finally, since $Y \to Y_0$ is
affine (as the base change of $X \to X_0$) we see that $Y$ has
an ample family of invertible modules by
Lemma \href{https://stacks.math.columbia.edu/tag/0GTQ}{lemma resolution property goes up affine (Tag 0GTQ)}.
This reduces us to the case discussed in the next paragraph.

\medskip\noindent
Assume $X$ is of finite type over $\mathbf{Z}$. In particular
$\dim(X) < \infty$. We will argue by induction on $\dim(X)$.
If $\dim(X) = 0$, then $X$ is affine and has the resolution property.
In general, there exists a dense open $U \subset X$ and a
$U$-admissible blowing up $X' \to X$ such that $X'$ has
an ample family of invertible modules, see
Lemma \href{https://stacks.math.columbia.edu/tag/0GTT}{lemma blow up ample family (Tag 0GTT)}.
Since $f : X' \to X$ is an isomorphism over $U$
we see that every point of $U$ lifts to a point of $X'$
with the same residue field.
Let $Z = X \setminus U$ with the reduced induced scheme structure.
Then $\dim(Z) < \dim(X)$ as $U$ is dense in $X$ (see above).
By induction we find a proper, completely decomposed
morphism $W \to Z$ such that $W$ has an ample family of invertible
modules. Then it follows that $Y = X' \amalg W \to X$ is the desired
morphism.
\end{proof}
\end{document}
